\section{Flatness}
\label{section:0.6}

\begin{env}[6.0]
\egalabel{0.6.0.0}{6.0.0}
The notion of flatness is due to J.-P.~Serre [16]; in what follows, we omit
the proofs of the results that are presented in the \emph{Algèbre
commutative} of N.~Bourbaki, to which we refer the reader. We
assume all rings to be commutative.\footnote{See the cited exposition by
N.~Bourbaki for the generalization of most of the results to the
noncommutative case.}

\oldpage[0\textsubscript{I}]{55}

If $M$, $N$ are two $A$-modules and $M'$ (resp. $N'$) is a submodule of $M$
(resp. $N$), we shall denote by $\operatorname{Im}(M'\otimes_A N')$ the submodule
of $M\otimes_A N$ that is the image of the canonical map
$M'\otimes_A N'\to M\otimes_A N$.
\end{env}

\subsection{Flat modules}
\label{subsection:0.6.1}

\begin{env}[6.1.1]
\label{0.6.1.1}
Let $M$ be an $A$-module. The following properties are equivalent:
\begin{enumerate}
  \item[$a)$] The functor $M\otimes_A N$ in $N$ is exact in the category
    of $A$-modules;
  \item[$b)$] $\operatorname{Tor}_i^A(M,N)=0$ for every $i>0$ and every
    $A$-module $N$;
  \item[$c)$] $\operatorname{Tor}_1^A(M,N)=0$ for every $A$-module $N$.
\end{enumerate}

When $M$ satisfies these conditions, $M$ is said to be a
\emph{flat $A$-module}. It is clear that every free $A$-module is flat.

For $M$ to be a flat $A$-module, it suffices that for every ideal
$\mathfrak{J}$ of $A$ \emph{of finite type}, the canonical map
$M\otimes_A\mathfrak{J}\to M\otimes_A A=M$ be \emph{injective}.
\end{env}

\begin{env}[6.1.2]\label{0.6.1.2}
\EGAReaderTarget{canon-ref.0.6.1.2-fr}{0.6.1.2-fr}
Every inductive limit of flat $A$-modules is a flat $A$-module. For
a direct sum $\bigoplus_{\lambda\in L}M_\lambda$ of $A$-modules to be
a flat $A$-module, it is necessary and sufficient that each of the $A$-modules
$M_\lambda$ be flat. In particular, every projective $A$-module is flat.

Let
\[
  0\longrightarrow M'\longrightarrow M\longrightarrow M''\longrightarrow 0
\]
be an exact sequence of $A$-modules such that $M''$ is \emph{flat}. Then, for
every $A$-module $N$, the sequence
\[
  0\longrightarrow M'\otimes N\longrightarrow M\otimes N
  \longrightarrow M''\otimes N\longrightarrow 0
\]
is exact. Moreover, for $M$ to be flat, it is then necessary and sufficient that
$M'$ be so (but it can happen that $M$ and $M'$ are flat without
$M''=M/M'$ being so).
\end{env}

\begin{env}[6.1.3]
\label{0.6.1.3}
Let $M$ be a flat $A$-module and $N$ an arbitrary $A$-module; for any two
submodules $N'$, $N''$ of $N$, we then have
\[
  \operatorname{Im}(M\otimes(N'+N''))
  =\operatorname{Im}(M\otimes N')+\operatorname{Im}(M\otimes N''),
\]
\[
  \operatorname{Im}(M\otimes(N'\cap N''))
  =\operatorname{Im}(M\otimes N')\cap\operatorname{Im}(M\otimes N'')
\]
(images taken in $M\otimes N$).
\end{env}

\begin{env}[6.1.4]
\label{0.6.1.4}
Let $M$, $N$ be two $A$-modules, $M'$ (resp. $N'$) a submodule of $M$
(resp. $N$), and suppose that one of the modules $M/M'$, $N/N'$ is flat.
Then we have
\[
  \operatorname{Im}(M'\otimes N')
  =\operatorname{Im}(M'\otimes N)\cap\operatorname{Im}(M\otimes N')
\]
(images in $M\otimes N$). In particular, if $\mathfrak{J}$ is an ideal
of $A$ and if $M/M'$ is flat, we have
$\mathfrak{J}M'=M'\cap\mathfrak{J}M$.
\end{env}

\subsection{Change of rings}
\label{subsection:0.6.2}

When an additive group $M$ is endowed with several module structures
relative to rings $A$, $B$, \ldots, instead of saying that $M$ is flat
as an $A$-module, $B$-module, \ldots, we shall sometimes also say that $M$
is \emph{$A$-flat}, \emph{$B$-flat}, \ldots

\begin{env}[6.2.1]
\label{0.6.2.1}
Let $A$ and $B$ be two rings, $M$ an $A$-module, $N$ an
$(A,B)$-bimodule. If $M$ is flat and $N$ is $B$-flat, then
$M\otimes_A N$ is $B$-flat. In particular, if $M$ and $N$ are two
flat $A$-modules, $M\otimes_A N$ is a flat $A$-module. If $B$ is an
$A$-algebra and $M$ is
\oldpage[0\textsubscript{I}]{56}
a flat $A$-module, the $B$-module $M_{(B)}=M\otimes_A B$ is flat. Finally, if
$B$ is an $A$-algebra that is flat as an $A$-module, and if $N$ is
a flat $B$-module, then $N$ is also $A$-flat.
\end{env}

\begin{env}[6.2.2]
\label{0.6.2.2}
Let $A$ be a ring and $B$ an $A$-algebra that is flat as an $A$-module.
Let $M$, $N$ be two $A$-modules such that $M$ admits a finite
presentation; then the canonical homomorphism
\[
  \operatorname{Hom}_A(M,N)\otimes_A B
  \longrightarrow
  \operatorname{Hom}_B(M\otimes_A B,N\otimes_A B)
  \tag{6.2.2.1}\label{0.6.2.2.1}
\]
(which sends $u\otimes b$ to the homomorphism
$m\otimes b'\longmapsto u(m)\otimes b'b$) is an isomorphism.
\end{env}

\begin{env}[6.2.3]
\label{0.6.2.3}
Let $(A_\lambda,\varphi_{\mu\lambda})$ be a filtered inductive system
of rings, and let $A=\varinjlim A_\lambda$. On the other hand, for each
$\lambda$, let $M_\lambda$ be an $A_\lambda$-module, and for $\lambda\leq\mu$ let
$\theta_{\mu\lambda}:M_\lambda\to M_\mu$ be a
$\varphi_{\mu\lambda}$-homomorphism, such that
$(M_\lambda,\theta_{\mu\lambda})$ is an inductive system;
$M=\varinjlim M_\lambda$ is then an $A$-module. That being so, if for every
$\lambda$, $M_\lambda$ is a \emph{flat} $A_\lambda$-module, then $M$ is
a \emph{flat} $A$-module. Indeed, let $\mathfrak{J}$ be an ideal \emph{of
finite type} of $A$; by the definition of the inductive limit, there exist an
index $\lambda$ and an ideal $\mathfrak{J}_\lambda$ of $A_\lambda$ such that
$\mathfrak{J}=\mathfrak{J}_\lambda A$. If we set
$\mathfrak{J}'_\mu=\mathfrak{J}_\lambda A_\mu$ for $\mu\geq\lambda$, we
also have $\mathfrak{J}=\varinjlim\mathfrak{J}'_\mu$ (where $\mu$ runs through the
indices $\geq\lambda$), whence (the functor $\varinjlim$ being exact and
commuting with the tensor product)
\[
  M\otimes_A\mathfrak{J}
  =\varinjlim(M_\mu\otimes_{A_\mu}\mathfrak{J}'_\mu)
  =\varinjlim\mathfrak{J}'_\mu M_\mu
  =\mathfrak{J}M.
\]
\end{env}

\subsection{Localization of flatness}
\label{subsection:0.6.3}

\begin{env}[6.3.1]
\label{0.6.3.1}
If $A$ is a ring and $S$ a multiplicative subset of $A$, $S^{-1}A$ is
a \emph{flat} $A$-module. Indeed, for every $A$-module $N$,
$N\otimes_A S^{-1}A$ is identified with $S^{-1}N$ (1.2.5), and we know (1.3.2)
that $S^{-1}N$ is an exact functor in $N$.

If now $M$ is a flat $A$-module,
$S^{-1}M=M\otimes_A S^{-1}A$ is a flat $S^{-1}A$-module (6.2.1), and hence
is also $A$-flat by virtue of the preceding and of (6.2.1). In particular,
if $P$ is an $S^{-1}A$-module, it can be considered as an $A$-module
isomorphic to $S^{-1}P$; for $P$ to be $A$-flat, it is necessary and sufficient
that it be $S^{-1}A$-flat.
\end{env}

\begin{env}[6.3.2]
\label{0.6.3.2}
Let $A$ be a ring, $B$ an $A$-algebra, and $T$ a multiplicative subset of
$B$. If $P$ is a $B$-module that is $A$-flat, $T^{-1}P$ is $A$-flat.
Indeed, for every $A$-module $N$, we have
\[
  (T^{-1}P)\otimes_A N
  =(T^{-1}B\otimes_B P)\otimes_A N
  =T^{-1}B\otimes_B(P\otimes_A N)
  =T^{-1}(P\otimes_A N)~;
\]
now, $T^{-1}(P\otimes_A N)$ is an exact functor in $N$, being the composite of the
two exact functors $P\otimes_A N$ (in $N$) and $T^{-1}Q$ (in $Q$). If $S$
is a multiplicative subset of $A$ whose image in $B$ is
\emph{contained in $T$}, $T^{-1}P$ is equal to $S^{-1}(T^{-1}P)$, and hence is
also $S^{-1}A$-flat by virtue of (6.3.1).
\end{env}

\begin{env}[6.3.3]
\label{0.6.3.3}
Let $\varphi:A\to B$ be a ring homomorphism and $M$ a $B$-module. The
following properties are equivalent:
\begin{enumerate}
  \item[$a)$] $M$ is a flat $A$-module.
  \item[$b)$] For every maximal ideal $\mathfrak{n}$ of $B$,
    $M_\mathfrak{n}$ is a flat $A$-module.
  \item[$c)$] For every maximal ideal $\mathfrak{n}$ of $B$, setting
    $\mathfrak{m}=\varphi^{-1}(\mathfrak{n})$, $M_\mathfrak{n}$ is a
    flat $A_\mathfrak{m}$-module.
\end{enumerate}

Indeed, since $M_\mathfrak{n}=(M_\mathfrak{n})_\mathfrak{m}$,
the equivalence of $b)$ and $c)$ follows from (6.3.1), and the fact that $a)$
implies $b)$ is a special case of (6.3.2). It remains to see that $b)$
implies $a)$,
\oldpage[0\textsubscript{I}]{57}
that is, that for every injective homomorphism $u:N'\to N$ of
$A$-modules, the homomorphism
$v=1\otimes u:M\otimes_A N'\to M\otimes_A N$ is injective. Now, $v$ is also
a homomorphism of $B$-modules, and we know that for it to be injective, it
suffices that for every maximal ideal $\mathfrak{n}$ of $B$,
\[
  v_\mathfrak{n}:(M\otimes_A N')_\mathfrak{n}
  \longrightarrow(M\otimes_A N)_\mathfrak{n}
\]
be injective. But since
\[
  (M\otimes_A N)_\mathfrak{n}
  =B_\mathfrak{n}\otimes_B(M\otimes_A N)
  =M_\mathfrak{n}\otimes_A N,
\]
$v_\mathfrak{n}$ is none other than the homomorphism
\[
  1\otimes u:M_\mathfrak{n}\otimes_A N'
  \longrightarrow M_\mathfrak{n}\otimes_A N,
\]
which is injective since $M_\mathfrak{n}$ is $A$-flat.

In particular (taking $B=A$), for an $A$-module $M$ to be flat, it is
necessary and sufficient that $M_\mathfrak{m}$ be $A_\mathfrak{m}$-flat for every
maximal ideal $\mathfrak{m}$ of $A$.
\end{env}

\begin{env}[6.3.4]
\label{0.6.3.4}
Let $M$ be an $A$-module; if $M$ is flat and if $f\in A$ is not a divisor
of $0$ in $A$, then $f$ annihilates no element $\neq0$ of $M$, because
the homomorphism $m\mapsto f\cdot m$ can be written $1\otimes u$, where $u$ is the
multiplication $a\mapsto f\cdot a$ in $A$ and $M$ is identified with
$M\otimes_A A$; if $u$ is injective, the same therefore holds for
$1\otimes u$, since $M$ is flat. In particular, if $A$ is
\emph{integral}, $M$ is \emph{torsion-free}.

Conversely, suppose $A$ integral and $M$ torsion-free, and suppose that for
every maximal ideal $\mathfrak{m}$ of $A$, $A_\mathfrak{m}$ is a
\emph{discrete valuation ring}; then $M$ is $A$-flat. Indeed, it suffices
(6.3.3) to prove that $M_\mathfrak{m}$ is $A_\mathfrak{m}$-flat, and
we can therefore assume that $A$ is already a discrete valuation ring. But
since $M$ is the inductive limit of its submodules of finite type, and these
are torsion-free, we can moreover restrict ourselves to the case where $M$ is of
finite type (6.1.2). In this case the proposition follows from the fact that $M$ is a
free $A$-module.

In particular, if $A$ is an \emph{integral} ring and
$\varphi:A\to B$ a ring homomorphism making $B$ an $A$-module that is
\emph{flat} and $\neq\{0\}$, then $\varphi$ is necessarily \emph{injective}.
Conversely, if $B$ is integral, $A$ a subring of $B$, and if for every
maximal ideal $\mathfrak{m}$ of $A$, $A_\mathfrak{m}$ is a discrete
valuation ring, then $B$ is $A$-flat.
\end{env}

\subsection{Faithfully flat modules}
\label{subsection:0.6.4}

\begin{env}[6.4.1]
\label{0.6.4.1}
For an $A$-module $M$, the following four properties are equivalent:
\begin{enumerate}
  \item[$a)$] For a sequence $N'\to N\to N''$ of $A$-modules to be exact,
    it is necessary and sufficient that the sequence
    $M\otimes N'\to M\otimes N\to M\otimes N''$ be exact;
  \item[$b)$] $M$ is flat and for every $A$-module $N$, the relation
    $M\otimes N=0$ implies $N=0$;
  \item[$c)$] $M$ is flat and for every homomorphism $v:N\to N'$ of
    $A$-modules, the relation $1_M\otimes v=0$ implies $v=0$, $1_M$ being
    the identity automorphism of $M$;
  \item[$d)$] $M$ is flat and for every maximal ideal $\mathfrak{m}$ of
    $A$, $\mathfrak{m}M\neq M$.
\end{enumerate}

When $M$ satisfies these conditions, $M$ is said to be a \emph{faithfully flat}
$A$-module; $M$ is then necessarily a \emph{faithful}
module. Moreover, if $u:N\to N'$ is a homomorphism of $A$-modules,
for $u$ to be injective (resp. surjective, bijective), it is necessary and sufficient
that $1\otimes u:M\otimes N\to M\otimes N'$ be so.
\end{env}

\oldpage[0\textsubscript{I}]{58}

\begin{env}[6.4.2]
\label{0.6.4.2}
A free module $\neq\{0\}$ is faithfully flat; the same holds for the
direct sum of a flat module and a faithfully flat module. If $S$ is
a multiplicative subset of $A$, $S^{-1}A$ is a faithfully flat $A$-module
only if $S$ consists of invertible elements (hence $S^{-1}A=A$).
\end{env}

\begin{env}[6.4.3]
\label{0.6.4.3}
Let
\[
  0\longrightarrow M'\longrightarrow M\longrightarrow M''\longrightarrow 0
\]
be an exact sequence of $A$-modules; if $M'$ and $M''$ are flat, and if one
of them is faithfully flat, then $M$ is faithfully flat.
\end{env}

\begin{env}[6.4.4]
\label{0.6.4.4}
Let $A$ and $B$ be two rings, $M$ an $A$-module, $N$ an
$(A,B)$-bimodule. If $M$ is faithfully flat and $N$ is a faithfully flat
$B$-module, then $M\otimes_A N$ is a faithfully flat $B$-module. In
particular, if $M$ and $N$ are two faithfully flat $A$-modules, the same
holds for $M\otimes_A N$. If $B$ is an $A$-algebra and $M$ is a
faithfully flat $A$-module, the $B$-module $M_{(B)}$ is faithfully flat.
\end{env}

\begin{env}[6.4.5]
\label{0.6.4.5}
If $M$ is a faithfully flat $A$-module and $S$ is a multiplicative
subset of $A$, $S^{-1}M$ is a faithfully flat $S^{-1}A$-module,
since $S^{-1}M=M\otimes_A(S^{-1}A)$ (6.4.4). Conversely, if for every
maximal ideal $\mathfrak{m}$ of $A$, $M_\mathfrak{m}$ is a faithfully flat
$A_\mathfrak{m}$-module, then $M$ is a faithfully flat $A$-module,
because $M$ is $A$-flat (6.3.3), and we have
\[
  M_\mathfrak{m}/\mathfrak{m}M_\mathfrak{m}
  =(M\otimes_A A_\mathfrak{m})
  \otimes_{A_\mathfrak{m}}
  (A_\mathfrak{m}/\mathfrak{m}A_\mathfrak{m})
  =M\otimes_A(A/\mathfrak{m})
  =M/\mathfrak{m}M,
\]
so the hypothesis implies that $M/\mathfrak{m}M\neq0$ for every maximal
ideal $\mathfrak{m}$ of $A$, which proves our assertion (6.4.1).
\end{env}

\subsection{Restriction of scalars}
\label{subsection:0.6.5}

\begin{env}[6.5.1]
\label{0.6.5.1}
Let $A$ be a ring and $\varphi:A\to B$ a ring homomorphism making
$B$ an $A$-algebra. Suppose there exists a $B$-module $N$ that is a
\emph{faithfully flat} $A$-module. Then, for every $A$-module $M$,
the homomorphism $x\mapsto 1\otimes x$ from $M$ to
$B\otimes_A M=M_{(B)}$ is \emph{injective}. In particular, $\varphi$ is
injective; for every ideal $\mathfrak{a}$ of $A$, we have
$\varphi^{-1}(\mathfrak{a}B)=\mathfrak{a}$; for every maximal (resp.
prime) ideal $\mathfrak{m}$ of $A$, there exists a maximal (resp. prime) ideal
$\mathfrak{n}$ of $B$ such that $\varphi^{-1}(\mathfrak{n})=\mathfrak{m}$.
\end{env}

\begin{env}[6.5.2]
\label{0.6.5.2}
When the condition of (6.5.1) is satisfied, we identify $A$ with a
subring of $B$ by means of $\varphi$, and more generally, for every $A$-module
$M$, we identify $M$ with a sub-$A$-module of $M_{(B)}$. Note that if
$B$ is then \emph{Noetherian}, so is $A$, because the map
$\mathfrak{a}\mapsto\mathfrak{a}B$ is an increasing injection of
the set of ideals of $A$ into the set of ideals of $B$;
the existence of an infinite strictly increasing sequence of ideals of $A$
would therefore imply the existence of an analogous sequence of ideals of $B$.
\end{env}

\subsection{Faithfully flat rings}
\label{subsection:0.6.6}

\begin{env}[6.6.1]
\label{0.6.6.1}
Let $\varphi:A\to B$ be a ring homomorphism making $B$ an
$A$-algebra. The following five properties are equivalent:
\begin{enumerate}
  \item[$a)$] $B$ is a faithfully flat $A$-module (in other words,
    $M_{(B)}$ is an exact and faithful functor in $M$).
  \item[$b)$] The homomorphism $\varphi$ is injective and the $A$-module
    $B/\varphi(A)$ is flat.

\oldpage[0\textsubscript{I}]{59}

  \item[$c)$] The $A$-module $B$ is flat (in other words, the functor
    $M_{(B)}$ is exact), and for every $A$-module $M$, the homomorphism
    $x\mapsto1\otimes x$ from $M$ to $M_{(B)}$ is injective.
  \item[$d)$] The $A$-module $B$ is flat and for every ideal $\mathfrak a$ of
    $A$, we have $\varphi^{-1}(\mathfrak aB)=\mathfrak a$.
  \item[$e)$] The $A$-module $B$ is flat and for every maximal ideal
    $\mathfrak m$ of $A$, there exists a maximal ideal $\mathfrak n$ of $B$ such
    that $\varphi^{-1}(\mathfrak n)=\mathfrak m$.
\end{enumerate}

When these conditions hold, we identify $A$ with a subring of $B$.
\end{env}

\begin{env}[6.6.2]
\label{0.6.6.2}
Let $A$ be a local ring, $\mathfrak m$ its maximal ideal, and $B$ an
$A$-algebra such that $\mathfrak mB\neq B$ (which is the case for example when
$B$ is a local ring and $A\to B$ a local homomorphism). If $B$ is a
flat $A$-module, $B$ is a faithfully flat $A$-module. Indeed, since
$\mathfrak mB\neq B$, there is a maximal ideal $\mathfrak n$ of $B$ containing
$\mathfrak mB$; since $\varphi^{-1}(\mathfrak n)\cap A$ contains
$\mathfrak m$ and does not contain $1$, we have
$\varphi^{-1}(\mathfrak n)=\mathfrak m$ and criterion $e)$ of (6.6.1)
applies. Under the stated conditions, we thus see that if $B$ is
Noetherian, so is $A$ (6.5.2).
\end{env}

\begin{env}[6.6.3]
\label{0.6.6.3}
Let $B$ be an $A$-algebra that is a faithfully flat $A$-module. For every
$A$-module $M$ and every sub-$A$-module $M'$ of $M$, we have (identifying $M$
with a sub-$A$-module of $M_{(B)}$) $M'=M\cap M'_{(B)}$. For $M$ to be a
flat (resp. faithfully flat) $A$-module, it is necessary and sufficient that $M_{(B)}$
be a flat (resp. faithfully flat) $B$-module.
\end{env}

\begin{env}[6.6.4]
\label{0.6.6.4}
Let $B$ be an $A$-algebra and $N$ a faithfully flat $B$-module. For $B$
to be a flat (resp. faithfully flat) $A$-module, it is necessary and sufficient that $N$
be so.

In particular, let $C$ be a $B$-algebra; if the ring $C$ is faithfully
flat over $B$ and $B$ is faithfully flat over $A$, then $C$ is faithfully flat
over $A$; if $C$ is faithfully flat over $B$ and over $A$, then $B$ is
faithfully flat over $A$.
\end{env}

\subsection{Flat morphisms of ringed spaces}
\label{subsection:0.6.7}

\begin{env}[6.7.1]
\label{0.6.7.1}
Let $f:X\to Y$ be a morphism of ringed spaces, and let $\mathcal F$ be an
$\mathcal O_X$-module. We say that $\mathcal F$ is $f$-flat (or $Y$-flat
when no confusion about $f$ is to be feared) at a point $x\in X$ if
$\mathcal F_x$ is a flat $\mathcal O_{f(x)}$-module; we say that
$\mathcal F$ is $f$-flat over $y\in Y$ if $\mathcal F$ is $f$-flat
at all the points $x\in f^{-1}(y)$; we say that $\mathcal F$ is $f$-flat
if $\mathcal F$ is $f$-flat at all the points of $X$. We say that the
morphism $f$ is flat at $x\in X$ (resp. flat over $y\in Y$, resp.
flat) if $\mathcal O_X$ is $f$-flat at $x$ (resp. $f$-flat over $y$,
resp. $f$-flat).
\end{env}

\begin{env}[6.7.2]
\label{0.6.7.2}
With the notation of (6.7.1), if $\mathcal F$ is $f$-flat at $x$, then for every
open neighbourhood $U$ of $y=f(x)$, the functor
$(f^*(\mathcal G)\otimes_{\mathcal O_X}\mathcal F)_x$ in $\mathcal G$ is
exact in the category of $(\mathcal O_Y|U)$-modules; indeed, this
stalk is canonically identified with
$\mathcal G_y\otimes_{\mathcal O_y}\mathcal F_x$, and our assertion follows
from the definition. In particular, if $f$ is a flat morphism, the functor
$f^*$ is exact in the category of $\mathcal O_Y$-modules.
\end{env}

\begin{env}[6.7.3]
\label{0.6.7.3}
Conversely, suppose the sheaf of rings $\mathcal O_Y$ is coherent, and
suppose that for every open neighbourhood $U$ of $y$, the functor
$(f^*(\mathcal G)\otimes_{\mathcal O_X}\mathcal F)_x$ is exact in
$\mathcal G$ in the category of coherent $(\mathcal O_Y|U)$-modules.
Then $\mathcal F$ is $f$-flat at $x$. Indeed, it suffices to prove that
for every ideal of finite type $\mathfrak J$ of $\mathcal O_y$,
the canonical homomorphism
$\mathfrak J\otimes_{\mathcal O_y}\mathcal F_x\to\mathcal F_x$ is injective
(6.1.1). Now, we know (5.3.8) that there then exist an open

\oldpage[0\textsubscript{I}]{60}

neighbourhood $U$ of $y$ and a coherent sheaf of ideals $\mathcal J$ of
$\mathcal O_Y|U$ such that $\mathcal J_y=\mathfrak J$, whence the conclusion.
\end{env}

\begin{env}[6.7.4]
\label{0.6.7.4}
The results of (6.1) on flat modules translate immediately into
propositions on sheaves that are $f$-flat at a point:

If $0\to\mathcal F'\to\mathcal F\to\mathcal F''\to0$ is an exact sequence
of $\mathcal O_X$-modules and if $\mathcal F''$ is $f$-flat at the point
$x\in X$, then, for every open neighbourhood $U$ of $y=f(x)$ and every
$(\mathcal O_Y|U)$-module $\mathcal G$, the sequence
\[
  0\longrightarrow
  (f^*(\mathcal G)\otimes_{\mathcal O_X}\mathcal F')_x
  \longrightarrow
  (f^*(\mathcal G)\otimes_{\mathcal O_X}\mathcal F)_x
  \longrightarrow
  (f^*(\mathcal G)\otimes_{\mathcal O_X}\mathcal F'')_x
  \longrightarrow0
\]
is exact. For $\mathcal F$ to be $f$-flat at $x$, it is then necessary and sufficient
that $\mathcal F'$ be so. There are analogous statements for the
corresponding notions of an $\mathcal O_X$-module $f$-flat over $y\in Y$, or
of an $f$-flat $\mathcal O_X$-module.
\end{env}

\begin{env}[6.7.5]
\label{0.6.7.5}
Let $f:X\to Y$, $g:Y\to Z$ be two morphisms of ringed spaces; let
$x\in X$, $y=f(x)$, and let $\mathcal F$ be an $\mathcal O_X$-module. If $\mathcal F$
is $f$-flat at the point $x$ and if the morphism $g$ is flat at the point $y$, then
$\mathcal F$ is $(g\circ f)$-flat at $x$ (6.2.1). In particular, if $f$ and
$g$ are flat morphisms, $g\circ f$ is flat.
\end{env}

\begin{env}[6.7.6]
\label{0.6.7.6}
Let $X$, $Y$ be two ringed spaces and $f:X\to Y$ a flat morphism. Then
the canonical homomorphism of bifunctors (4.4.6)
\[
  f^*(\mathcal{H}om_{\mathcal O_Y}(\mathcal F,\mathcal G))
  \longrightarrow
  \mathcal{H}om_{\mathcal O_X}(f^*(\mathcal F),f^*(\mathcal G))
  \tag{6.7.6.1}
  \label{0.6.7.6.1}
\]
is an isomorphism when $\mathcal F$ admits a finite presentation
(5.2.5).

Indeed, the question being local, we can assume that there exists an exact
sequence $\mathcal O_Y^m\to\mathcal O_Y^n\to\mathcal F\to0$. Now, both
sides of (6.7.6.1) are left exact functors in $\mathcal F$ by
virtue of the hypothesis on $f$; we are then reduced to proving the proposition
when $\mathcal F=\mathcal O_Y$, a case in which it is trivial.
\end{env}

\begin{env}[6.7.8]
\label{0.6.7.8}
A morphism $f:X\to Y$ of ringed spaces is said to be faithfully flat if
$f$ is surjective and if, for every $x\in X$, $\mathcal O_x$ is a
faithfully flat $\mathcal O_{f(x)}$-module. When $X$ and $Y$ are
locally ringed spaces (5.5.1), this amounts to saying that the
morphism $f$ is surjective and flat (6.6.2). When $f$ is faithfully flat,
$f^*$ is an exact and faithful functor in the category of
$\mathcal O_Y$-modules (6.6.1, $a)$), and for an $\mathcal O_Y$-module
$\mathcal G$ to be $Y$-flat, it is necessary and sufficient that $f^*(\mathcal G)$ be
so (6.6.3).
\end{env}

